\documentclass[11pt,letterpaper]{article}
\usepackage[T1]{fontenc}
\usepackage{amsmath,amssymb,amsthm,mathtools,booktabs,array,microtype}
\usepackage[margin=1in]{geometry}
\usepackage[colorlinks=true,linkcolor=blue,citecolor=blue,urlcolor=blue]{hyperref}
\hypersetup{pdftitle={Pure fibres of the four-site matching map},pdfauthor={Ignoramus}}
\allowdisplaybreaks
\newcommand{\C}{\mathbb C}
\newcommand{\A}{\mathbb A}
\newcommand{\Gm}{\mathbb G_m}
\newcommand{\rk}{\operatorname{rank}}
\newcommand{\im}{\operatorname{im}}
\newcommand{\Span}{\operatorname{span}}
\newcommand{\Ftau}{F_\tau}
\newcommand{\Fact}{F_\tau^{\mathrm{nz}}}
\newcommand{\Ddec}{D_\tau}
\newtheorem{theorem}{Theorem}[section]
\newtheorem{proposition}[theorem]{Proposition}
\newtheorem{lemma}[theorem]{Lemma}
\newtheorem{corollary}[theorem]{Corollary}
\theoremstyle{definition}
\newtheorem{definition}[theorem]{Definition}
\newtheorem{remark}[theorem]{Remark}
\title{Pure fibres of the four-site matching map}
\author{Ignoramus}
\date{11 November 2026}
\begin{document}
\maketitle
\begin{abstract}
The four-site matching map sends six pair tensors on four three-dimensional
complex vector spaces to the sum of their products over the three perfect
matchings. Over a general output its fibre is known to be a single orbit of
a three-dimensional torus. We study the fibre over a pure tensor, that is,
a nonzero product of four vectors. We give complete parametrizations of
this fibre, organized by nine support and rank patterns after quotienting
each local space by its target line. The fibre has dimension \(19\) and
exactly twelve irreducible components of that dimension, each isomorphic to
\(\A^{18}\times\Gm\). On the open subset where all six pair tensors are
nonzero, the dimension is \(15\), again with twelve components of maximal
dimension. The corresponding space of three ordered nonzero summands, one
decomposable across each of the three pair partitions, has dimension \(12\).
The proof combines a projected-support obstruction, elementary tensor
syzygies, and an incidence calculation for isotropic planes. The
parametrizations allow zero blocks and drops of whole rank;
lower-dimensional components are not classified.
\end{abstract}

\section{The fibre and the main results}

Let \(V_1,\ldots,V_4\) be three-dimensional complex vector spaces. Put
\(U_{ij}=V_i\otimes V_j\) for \(i<j\), and consider the polynomial map
\[
 \Phi:\prod_{i<j}U_{ij}\longrightarrow
 V_1\otimes V_2\otimes V_3\otimes V_4,\qquad
 \Phi(W)=W_{12}W_{34}+W_{13}W_{24}+W_{14}W_{23}.
\]
A product in this formula is the tensor product on disjoint sites, restored
to the order \(1,2,3,4\), with no sign. Thus the three terms correspond to
the three perfect matchings of the complete graph on four labelled vertices.
The map \(\Phi\) is the matching tensor of that graph in the sense
of~\cite{ign-geometry}, here with all local dimensions equal to three.
Matching tensors arise from the description of optical quantum experiments
by graphs~\cite{kgz}; see the introduction of~\cite{ign-geometry}. No
question from that subject is decided here. For three colours and no
internal vertices, the same map gives the full-sector coefficients in the
model of \cite[Section~2]{ign-boundary}, of which no result is used here.

The closure of the image of \(\Phi\) is the join of the three cones of
tensors decomposable across \(12|34\), \(13|24\) and \(14|23\)
\cite[Proposition~7.1 and its proof]{ign-geometry}, and over a general
point of that closure the whole fibre is a single orbit of the
three-dimensional torus of vertex rescalings
\cite[Theorem~4.5]{ign-geometry}. That paper lists a description of the
exceptional outputs among the things it does not prove. Here an output is
called \emph{exceptional} if its fibre is not a single such orbit. Call a
tensor in \(V_1\otimes V_2\otimes V_3\otimes V_4\) \emph{pure} if it is a
nonzero product of four vectors, one from each \(V_i\), and call the fibre
of \(\Phi\) over a pure tensor a \emph{pure fibre}. A pure tensor lies in
all three cones and is exceptional: its fibre has dimension \(19\), not
three (Theorem~\ref{main:geometry}).

Fix nonzero \(c_i\in V_i\) and \(\tau\in\C^*\). We study the reduced affine
algebraic set
\[
 \Ftau=\{W:\Phi(W)=\tau c_1c_2c_3c_4\}.
\]
No quotient by changes of basis is taken. Let
\[
 \Fact=\{W\in\Ftau:W_{ij}\ne0\text{ for all }i<j\}.
\]
This is an open subvariety. Its defining nonvanishing condition is equivalent
to requiring all three matching products to be nonzero.

The same question can be asked of the three products alone: in which ways
is a pure tensor a sum of three nonzero tensors, one decomposable across
each of the three pair partitions? Let \(\mathcal S_{ij|kl}^{\times}\) be
the set of nonzero tensors of the form \(XY\) with \(X\in U_{ij}\) and
\(Y\in U_{kl}\), the product being formed as in the definition of \(\Phi\).
Define
\[
 \Ddec=\left\{(A,B,C)\in
 \mathcal S_{12|34}^{\times}\times
 \mathcal S_{13|24}^{\times}\times
 \mathcal S_{14|23}^{\times}:
 A+B+C=\tau c_1c_2c_3c_4\right\}.
\]
The summands are ordered and their three partitions are prescribed. They need
not be pure.

\begin{theorem}\label{main:geometry}
With the preceding hypotheses:
\begin{enumerate}
\item \(\dim\Ftau=19\). Exactly twelve irreducible components have dimension
\(19\), and each is isomorphic to \(\A^{18}\times\Gm\).
\item \(\dim\Fact=15\). Exactly twelve irreducible components of \(\Fact\)
have dimension \(15\).
\item The matching-product map \(\Fact\to\Ddec\) is a Zariski locally trivial
bundle with fibre \((\Gm)^3\). Consequently \(\dim\Ddec=12\), and exactly
twelve irreducible components of \(\Ddec\) have dimension \(12\).
\end{enumerate}
These assertions specify the components of maximal dimension, not the
complete irreducible decompositions of these varieties.
\end{theorem}

The large components of the full fibre are easy to exhibit. Choose a vertex
\(i\), two of its neighbours \(j,k\), and write \(z\) for the remaining vertex.
Set
\begin{equation}\label{main:large}
 W_{jz}=W_{kz}=0,\qquad
 W_{iz}=q c_ic_z,\qquad W_{jk}=(\tau/q)c_jc_k,
 \qquad q\ne0,
\end{equation}
and leave \(W_{ij},W_{ik}\) arbitrary. The two free blocks do not contribute
to the output because their complementary blocks vanish. The first assertion
of the theorem includes the global upper bound and the fact that these are
all components of that dimension. The second and third assertions concern
cancellation among three nonzero products, where no block is free in this
way. There the components of dimension fifteen are the closures in
\(\Fact\) of the strata whose projected support, defined below, is a
triangle with ranks \((2,1,1)\); the
chart~\eqref{geom:triangle-chart} parametrizes these strata, and
Section~\ref{examples:section} gives an explicit point.

Here is the classification used to prove the theorem. Let
\(\pi_i:V_i\to V_i/\C c_i\), and write
\(\overline W_{ij}=(\pi_i\otimes\pi_j)W_{ij}\).
The \emph{projected support graph} has edge \(ij\) precisely when
\(\overline W_{ij}\ne0\). A zero projected block may have a nonzero whole
block.

A block \(W_{ij}\) is \emph{pure} if it is a nonzero multiple of
\(c_ic_j\). The \emph{star} at a vertex is the set of the three edges at
it, and that vertex is its \emph{centre}. A \emph{wedge} is a projected
support consisting of two adjacent edges, \emph{centred} at their common
vertex. The other two vertices of a wedge, and the three vertices other
than the centre of a star, are its \emph{leaves}. If the projected support
is a triangle, the blocks joining its vertices to the fourth vertex are its
\emph{spokes}; if its projected ranks are \((2,1,1)\), the edge of rank two
is the \emph{high} edge and the other two are the \emph{low} edges. A point
of \(\Ftau\) is \emph{active} if it lies in \(\Fact\); an active stratum is
the intersection of a stratum with \(\Fact\), and an active component is an
irreducible component of \(\Fact\). The exceptional site in
Lemma~\ref{atlas:crossing} only labels a row of its table; it is unrelated
to exceptional outputs.

\begin{theorem}\label{main:atlas}
Every point of \(\Ftau\), up to simultaneous relabelling of the four sites,
has exactly one of the nine projected support and rank patterns in
Table~\ref{main:table}. Every pattern occurs. Section~\ref{atlas:section}
gives complete whole-tensor parametrizations of all nine types.
The dimension in the table is that of a stratum with its labelled support
and rank assignment fixed.
\end{theorem}

\begin{table}[ht]
\centering
\begin{tabular}{@{}llrr@{}}
\toprule
Projected support & Nonzero projected ranks & Dimension & Nonzero-block part\\
\midrule
Empty & --- & 11 & \(\le11\)\\
One edge & \(1\) & 14 & 12\\
One edge & \(2\) & 15 & 13\\
Adjacent edges & \(1,1\) & 17 & 13\\
Adjacent edges & \(2,1\) & 18 & empty\\
Adjacent edges & \(2,2\) & 19 & empty\\
Triangle & \(1,1,1\) & 14 & 14\\
Triangle & \(2,1,1\) & 15 & 15\\
Four-cycle & \(1,1,1,1\) & 12 & 12\\
\bottomrule
\end{tabular}
\caption{Complete projected types. The last column concerns the intersection
with \(\Fact\); only an upper bound is needed in its first row.}
\label{main:table}
\end{table}

The classification is a parametrization by finitely many types, with
redundant parameters where explicitly indicated. It is not a classification
of orbits under local linear groups. In particular the empty-support
parametrization uses isotropic subspaces, rather than a finite list of
coordinate representatives.

\subsection*{Relation to earlier work}
The problem is a nonminimal decomposition problem at an exceptional output,
for three different paired Segre cones.
Generic uniqueness results for joins, such as
Ballico's Proposition~8~\cite{ballico}, assume a general point of the relevant
join and its prescribed join rank. The target here belongs to each
constituent Segre cone. Nonminimal decomposition results of
Lovitz and Petrov~\cite[Section~6.2]{lovitzpetrov} concern product tensors
relative to one fixed factorization; the three factorizations above differ.
Neither statement determines this fibre.

The ingredients of our proof have familiar counterparts. The intersection of
two crossing partition Segre varieties is the Segre variety for their common
refinement; see the proof of Theorem~6.1 in Section~6 of Gesmundo, Han and
Lovitz~\cite{ghl}, in the numbering of arXiv:2407.16767v2. This accounts
for the rank-one factors in the four-cycle case. Zero-slice decompositions,
as studied by Karam~\cite{karam} (Proposition~3.2 in the numbering of
arXiv:2308.07101v1), provide a broader setting for the elementary syzygies
used after a triangle reduction. The two-term case used repeatedly below,
in which a relation \(xA+yB=0\) between nonzero tensors \(A,B\), with
nonzero vectors \(x,y\) at two different sites, forces \(A=hy\) and
\(B=-hx\) for a common factor \(h\), is the two-leaf case of
\cite[Theorem~2.1]{ign-relations}; see also Corollary~9.2 there.
The proof of Lemma~\ref{support:degree} contracts with a covector chosen
outside finitely many kernels. The same generic-contraction and
finite-union device appears, in another setting, in Section~3 of the
public provisional note~\cite{axis}; no proof here depends on that note.
The quadratic-form calculations below use ordinary isotropic-subspace
geometry~\cite{ekm}. The surjectivity of the bundle in
Theorem~\ref{main:geometry}(3) is the four-vertex construction in the
proof of~\cite[Proposition~7.1]{ign-geometry}.

Our subject is the complete fibre over a pure target in this fixed format,
including its whole-tensor lifts and its sharp dimension bounds. The result
does not assert a general theorem on partition rank~\cite{naslund} or an
exclusion for larger matching systems.
Bounded searches by AI agents found no earlier statement of the results
of this paper; this is not proof that they are new.

\section{Exhaustion of the projected supports}

All tensor products in the proofs retain the actual given blocks. We may
choose splittings \(V_i=\C c_i\oplus H_i\), with \(\dim H_i=2\), when
coordinates are useful.

\begin{lemma}\label{support:degree}
The projected support graph of a point of \(\Ftau\) has maximum degree at
most two.
\end{lemma}
\begin{proof}
Suppose that all three projections incident to a vertex \(i\) are nonzero.
Their contractions are three nonzero linear maps from
\(\operatorname{Ann}(c_i)\) to the corresponding quotient spaces. A finite
union of proper linear subspaces cannot cover a complex vector space.
Choose \(\alpha\in\operatorname{Ann}(c_i)\) outside their three kernels.

Write the other vertices, the leaves of the star at \(i\), as \(j,k,z\).
The whole contractions \(f_j,f_k,f_z\)
of the incident blocks by \(\alpha\) have nonzero quotient images, so each
\(f_r\) is independent of \(c_r\). Choose fixed dual functionals
\(\delta_r,\epsilon_r\) satisfying
\[
 \delta_r(c_r)=1,\quad \delta_r(f_r)=0,\qquad
 \epsilon_r(c_r)=0,\quad \epsilon_r(f_r)=1.
\]
Contract the leaf slots by
\(\delta_j+u\epsilon_j,\delta_k+v\epsilon_k,\delta_z+w\epsilon_z\).
Keeping the centre slot whole gives an identity
\[
 A(u)\phi(v,w)+B(v)\psi(u,w)+C(w)\chi(u,v)=\tau c_i.
\]
The scalar factors have the displayed absent-variable dependencies because
their blocks do not meet the missing leaf. Applying \(\alpha\) yields
\[
 u\phi(v,w)+v\psi(u,w)+w\chi(u,v)=0.
\]
Evaluation at \((1,0,0),(0,1,0),(0,0,1)\) gives
\(\phi(0,0)=\psi(0,0)=\chi(0,0)=0\). The whole identity at the origin would then say
\(0=\tau c_i\), a contradiction. This proof allows zero leaf blocks and
arbitrary components along the target lines.
\end{proof}

\begin{lemma}\label{support:triangle}
If the projected support is a triangle, at least two of its edges have
rank one. More precisely, if \(z\) is the isolated projected vertex, there
is a unique triangle vertex \(i\) whose whole spoke to \(z\) is a nonzero
multiple of \(c_ic_z\). The two triangle edges incident to \(i\) have
rank-one projections.
\end{lemma}
\begin{proof}
Write the triangle vertices as \(l,m,t\), with
\[
 L=W_{lm},\quad P=W_{lt},\quad R=W_{mt},\qquad
 Q=W_{lz},\quad S=W_{mz},\quad K=W_{tz}.
\]
The three triangle projections are nonzero; the other three vanish.
Projecting at \(l,m,z\) in the full equation isolates
\(\overline L(\mathrm{id}_t\otimes\pi_z)K=0\). Thus
\(K=p_t c_z\). The other two triple projections give
\(Q=p_l c_z,S=p_m c_z\). Cancelling the fixed nonzero factor \(c_z\),
\[
 Lp_t+Pp_m+Rp_l=\tau c_lc_mc_t.
\]
Quotienting each triangle space by its actual line \(\Span(p_r)\), with
\(\Span(0)=0\), kills the left side. Some \(c_i\) must therefore belong
to \(\Span(p_i)\). Write \(p_i=\lambda c_i\), where \(\lambda\ne0\).

Let \(j,k\) be the other triangle vertices and put
\(A=W_{ij},B=W_{ik},C=W_{jk}\). Then
\[
 Ap_k+Bp_j+\lambda c_iC=\tau c_ic_jc_k.
\]
If \(\bar p_j=0\), projection at \(i,j\) forces \(p_k=0\), and projection at
\(i,k\) then forces \(p_j=0\). The equation would imply \(\bar C=0\),
a contradiction. Similarly \(\bar p_k\ne0\). This also proves uniqueness
of the pure spoke.
Indeed a pure spoke at \(j\) or \(k\) would mean \(p_j\in\C c_j\) or
\(p_k\in\C c_k\), that is, \(\bar p_j=0\) or \(\bar p_k=0\).

Project only at \(i\), writing \(A',B'\) for the resulting tensors.
The identity \(A'p_k+B'p_j=0\), with \(p_j,p_k\ne0\), implies
\(A'=hp_j,B'=-hp_k\) for one \(h\ne0\): apply the quotient by \(\Span(p_j)\)
at \(j\), then cancel the nonzero factors. Choose a whole lift \(u\) of \(h\).
The complete whole blocks have the form
\begin{equation}\label{support:triangle-form}
 A=up_j+c_iv,\qquad B=-up_k+c_iw,\qquad
 C=\lambda^{-1}(\tau c_jc_k-vp_k-p_jw).
\end{equation}
In particular \(\bar A=h\bar p_j\) and \(\bar B=-h\bar p_k\) have rank one.
No initial rank-one assumption was used.
\end{proof}

\begin{proof}[Proof of the support assertion in Theorem~\ref{main:atlas}]
A simple graph on four vertices of maximum degree two is empty, one edge,
two adjacent edges, two disjoint edges, a three-edge path, a triangle,
or a four-cycle. The disjoint-edge and three-edge-path graphs each have
exactly one supported perfect matching. Fourfold projection of \(\Phi(W)\)
would then be one nonzero product, whereas the projected target is zero.

For a four-cycle, take its missing edges to be \(lm,tz\).
The fourfold equation reads
\(\overline P\,\overline S=-\overline Q\,\overline R\ne0\).
Across \(lt|mz\) the left side has rank one and the right side has rank
\(\rk\overline Q\,\rk\overline R\), since its matrix is a reordered
Kronecker product. Both ranks on the right are one. The other crossing
cut gives the same conclusion for \(P,S\).
Lemma~\ref{support:triangle} leaves only ranks \(111\) and \(211\) for
a triangle. All remaining nonzero ranks are one or two because the local
quotients have dimension two. This gives the nine types.
The complete lift charts below establish sufficiency and nonemptiness.
For instance, in \eqref{main:large}, free blocks \(W_{ij},W_{ik}\) of
projected ranks \((0,0)\), \((1,0)\), \((2,0)\), \((1,1)\), \((2,1)\) and
\((2,2)\) give the first six types of Table~\ref{main:table} in order;
\eqref{atlas:triangle-form} with \(v=q\) and \(w=0\) gives a triangle of
ranks \((1,1,1)\); the point of Section~\ref{examples:section} is a
triangle of ranks \((2,1,1)\); and \eqref{atlas:cycle-form} with
\(\gamma=1\) and all \(u_i\in H_i\setminus\{0\}\) gives a four-cycle.
\end{proof}

\section{Whole-tensor parametrizations}\label{atlas:section}

We give the whole tensors in each of the nine types.  In this section
\((l,m,t,z)\) is an ordering of the four sites, and
\[
 L=W_{lm},\quad P=W_{lt},\quad Q=W_{lz},\quad
 R=W_{mt},\quad S=W_{mz},\quad K=W_{tz}.
\]
Juxtaposition denotes tensor product, restored without a sign to the indicated
site order.  The equation to be solved is
\begin{equation}\label{atlas:equation}
 LK+PS+QR=\tau c_lc_mc_tc_z,\qquad \tau\ne0.
\end{equation}
Write \(V_i=\C c_i\oplus H_i\), using the fixed complements to identify
quotients with \(H_i\).  A bar on a vector is its quotient image; on an edge
tensor it means projection at both endpoints.  In particular,
\begin{equation}\label{atlas:zero-kernel}
 \overline{W}_{ij}=0
 \quad\Longleftrightarrow\quad
 W_{ij}\in c_i\otimes V_j+V_i\otimes c_j.
\end{equation}
A tensor described as free with prescribed projected rank may have arbitrary
entries involving a \(c\) factor, including entries making its whole matrix
singular.  All scalar and vector parameters below are complex.

\subsection{Empty projected support}

Let \(E=\C^6\), with coordinates indexed by the six unordered pairs of sites,
and put
\[
 q(x)=x_{12}x_{34}+x_{13}x_{24}+x_{14}x_{23},\qquad
 b(x,y)=q(x+y)-q(x)-q(y).
\]
Let \(S_i\subset E\) be the coordinate subspace on the three edges incident
to \(i\).

\begin{proposition}\label{atlas:empty}
The solutions with all six double projections zero are parametrized as follows.
Choose \(a\in E\) with \(q(a)=\tau\), a totally isotropic subspace
\(U\subset a^\perp\) for \(b\), and linear maps
\[
 B_i:H_i^*\longrightarrow U\cap S_i\qquad(1\le i\le4).
\]
Necessarily \(\dim U\le2\).  Define \(u_{ij}\in H_i\) by
\((B_i(x))_{ij}=x(u_{ij})\), and set
\begin{equation}\label{atlas:empty-form}
 W_{ij}=a_{ij}c_ic_j+u_{ij}c_j+c_iu_{ji}.
\end{equation}
These formulas give every solution of this type.  The data \(a,B_i\) are
unique; the choice of \(U\) becomes unique if
\(U=\sum_i\im B_i\) is required.
\end{proposition}

\begin{proof}
By \eqref{atlas:zero-kernel}, every such tuple has the unique form
\eqref{atlas:empty-form}.  Evaluate the tensor at covectors taking \(c_i\)
to \(1\), and denote their restrictions to \(H_i\) by \(x_i\).
The six edge evaluations form the vector
\(a+\sum_i B_i(x_i)\).  Thus the target identity is equivalent to
\[
 q\left(a+\sum_iB_i(x_i)\right)=\tau
 \quad\text{for all }x_i\in H_i^*.
\]
This equivalence loses no coefficients: the expression is multiaffine in
the four groups of variables, and homogenizing each group recovers the
original four-tensor.  Comparison of constant, linear and quadratic terms,
followed by polarization, gives precisely
\[
 q(a)=\tau,\qquad b(a,U)=0,\qquad b(U,U)=0,
 \qquad U=\sum_i\im B_i.
\]
The quadratic terms give \(b(\im B_i,\im B_j)=0\) for \(i\ne j\); for
\(i=j\) this holds because each \(S_i\) is totally isotropic, no two
edges at \(i\) being complementary.
Conversely these conditions make the displayed polynomial constant.
The form \(b\) is nondegenerate, since its matrix has three invertible
complementary-edge blocks.  As \(b(a,a)=2\tau\ne0\), the space \(a^\perp\)
is nondegenerate of dimension five.  For an isotropic \(U\) in that space,
\(U\subset U^\perp\), so \(2\dim U\le5\).
\end{proof}

\subsection{A single projected edge}

We first record the crossing decomposition needed at a zero opposite edge.

\begin{lemma}\label{atlas:crossing}
Suppose \(\overline P=\overline Q=\overline R=\overline S=0\) and
\(PS+QR=\tau c_lc_mc_tc_z\).  The following list is exhaustive.
\begin{enumerate}
\item If \(PS=0\), then for some \(a\ne0\),
\(Q=a c_lc_z\), \(R=(\tau/a)c_mc_t\), and \(P=0\) or \(S=0\).
The unused partner is any zero-double-projection tensor.
\item If \(QR=0\), then for some \(a\ne0\),
\(P=a c_lc_t\), \(S=(\tau/a)c_mc_z\), and \(Q=0\) or \(R=0\).
The unused partner is any zero-double-projection tensor.
\item If both products are nonzero, all four edges have whole rank one.
Choose one exceptional site and one row of the following table:
\[
\begin{array}{c|cccc|c}
 &P&S&Q&R&\text{condition}\\ \hline
 l&u c_t&s c_mc_z&v c_z&r c_mc_t&su+rv=\tau c_l\\
 m&p c_lc_t&u c_z&q c_lc_z&v c_t&pu+qv=\tau c_m\\
 t&c_lu&s c_mc_z&q c_lc_z&c_mv&su+qv=\tau c_t\\
 z&p c_lc_t&c_mu&c_lv&r c_mc_t&pu+rv=\tau c_z
\end{array}
\]
All named scalar parameters and both vectors \(u,v\) are nonzero.
In each row the displayed equation explicitly determines \(v\) from \(u\)
and the two scalars; the only additional restriction is \(v\ne0\).
The choices of exceptional site may overlap.
\end{enumerate}
Every listed choice is a solution.
\end{lemma}

\begin{proof}
If one product vanishes, the other equals the nonzero pure tensor.
Factorization across its own edge bipartition gives the first two cases;
a disjoint tensor product vanishes exactly when a whole factor vanishes.

Suppose both products are nonzero.  Every edge has rank at most two by
\eqref{atlas:zero-kernel}.  If \(P\) had rank two, flatten
\(PS=\tau c_lc_mc_tc_z-QR\) across \((l,z)\mid(m,t)\).
Its rank is \(\rk(P)\rk(S)\), whereas the right side has rank at most two.
Thus \(S=s_ms_z\) has rank one and the flattening has rank two.
Its left and right spaces are respectively
\[
 \im_l(P)\otimes\C s_z=\Span(c_lc_z,Q),\qquad
 \C s_m\otimes\im_t(P)=\Span(c_mc_t,R).
\]
It follows that \(S=s c_mc_z\), \(Q=q c_z\), and \(R=c_mv\), with
\(s\ne0\).  Hence \(sP=\tau c_lc_t-qv\).
Rank two of \(P\) forces \(q\notin\C c_l\) and \(v\notin\C c_t\),
contradicting \(\overline P=0\).  The same argument, with the other crossing
cut, applies to each of \(Q,R,S\).  Thus all four have rank one.

Now \(PS\) and \(QR\) are nonzero simple four-tensors whose sum is nonzero
and simple.  Their factor lines can differ at at most one site: otherwise
flattening at one differing site gives two independent factors on each side
and therefore rank two.  At every other site the common line must be the
target line \(\C c_i\).  Absorbing the scalar factors gives exactly the four
rows.  Substitution proves their converses, including the case in which all
four factor lines agree.
\end{proof}

\begin{proposition}\label{atlas:single}
Fix \(r\in\{1,2\}\).  Suppose only \(\overline L\) is nonzero and its rank
is \(r\).  Then \(K=\kappa c_tc_z\), and all solutions are the following.
\begin{enumerate}
\item If \(\kappa=0\), take \(L\) free with projected rank \(r\), and take
\(P,Q,R,S\) from Lemma~\ref{atlas:crossing}.
\item If \(\kappa\ne0\) and both crossing products are nonzero, choose
nonzero \(p,q\in V_l\) and \(r_0,s\in V_m\), and put
\begin{equation}\label{atlas:single-live}
 \begin{gathered}
 P=p c_t,\quad Q=q c_z,\quad R=r_0c_t,\quad S=s c_z,\\
 L=\kappa^{-1}(\tau c_lc_m-ps-qr_0).
 \end{gathered}
\end{equation}
Require the matrix \(M=\overline p\,\overline s+
\overline q\,\overline{r_0}\) to have rank \(r\).
Explicitly, in any quotient bases, this means
\[
\begin{cases}
 M\ne0,\quad
 \det[\overline p,\overline q]\det[\overline s,\overline{r_0}]=0,
 &r=1,\\
 \det[\overline p,\overline q]\det[\overline s,\overline{r_0}]\ne0,
 &r=2.
\end{cases}
\]
\item If \(\kappa\ne0\) and \(QR=0\), necessarily \(r=1\).
Choose \(p\in V_l,s\in V_m\) with nonzero quotients, and set
\[
 P=p c_t,\quad S=s c_z,\quad
 L=\kappa^{-1}(\tau c_lc_m-ps).
\]
Choose any zero-double-projection \(Q,R\) with \(Q=0\) or \(R=0\).
\item If \(\kappa\ne0\) and \(PS=0\), necessarily \(r=1\).
Choose \(q\in V_l,r_0\in V_m\) with nonzero quotients, and set
\[
 Q=q c_z,\quad R=r_0c_t,\quad
 L=\kappa^{-1}(\tau c_lc_m-qr_0).
\]
Choose any zero-double-projection \(P,S\) with \(P=0\) or \(S=0\).
\end{enumerate}
The branches exhaust all zero and nonzero crossing products.
\end{proposition}

\begin{proof}
Project \eqref{atlas:equation} at \(l,m,t\), and then at \(l,m,z\).
The crossing terms vanish in each case, so the nonzero \(\overline L\)
forces both partial quotients of \(K\) to vanish.  Hence \(K=\kappa c_tc_z\).
Write the other edges uniquely as
\[
 P=p c_t+c_la,\quad Q=q c_z+c_lb,\quad
 R=r_0c_t+c_md,\quad S=s c_z+c_me,
\]
where \(a,d\in H_t\), \(b,e\in H_z\), and the other vectors are whole.
Separating the four \((t,z)\) sectors gives the equivalent full system
\begin{equation}\label{atlas:single-sectors}
\begin{aligned}
 \kappa L+ps+qr_0&=\tau c_lc_m,\\
 c_lsa+qc_md&=0,&
 pc_me+c_lr_0b&=0,& ae+db&=0.
\end{aligned}
\end{equation}
When \(\kappa=0\), \(L\) drops out and Lemma~\ref{atlas:crossing} applies.

Let \(\kappa\ne0\).  The double projection of the first equation shows
that at least one of \(\overline p\,\overline s\) and
\(\overline q\,\overline{r_0}\) is nonzero.  If the first is nonzero,
projecting the second equation at \(m\) gives \(a=0\), and projecting the
third at \(l\) gives \(e=0\).  The remaining conditions are
\(qd=0\), \(r_0b=0\), and \(db=0\).  If \(Q,R\ne0\), these force
\(b=d=0\): for example \(d\ne0\) would imply \(q=b=0\), hence \(Q=0\).
If \(QR=0\), the unused partner is unrestricted within
\eqref{atlas:zero-kernel}.  Interchanging the two crossing products proves
the corresponding assertion when the second projected summand is nonzero.
Both whole products cannot vanish, since that would make \(\overline L=0\).
This proves all the stated forms.  The determinant formula follows by
writing \(M=[\overline p,\overline q][\overline s,\overline{r_0}]^T\).
A single nonzero projected outer product has rank one, explaining the last
two branches.  Finally the displayed formulas solve the entire system
\eqref{atlas:single-sectors}, which proves every converse.
\end{proof}

\subsection{Two adjacent projected edges}

Here \(\overline L,\overline P\ne0\) and
\(\overline Q=\overline R=\overline S=\overline K=0\).

\begin{lemma}\label{atlas:wedge-reduction}
Every solution in this support has
\begin{equation}\label{atlas:wedge-reduced}
 Q=q c_z,\quad S=s c_z,\quad K=k c_z,\qquad
 Lk+Ps+qR=\tau c_lc_mc_t,
\end{equation}
where \(q\ne0\) and \(R\ne0\).  Either both \(\overline s,\overline k\)
vanish or both are nonzero.
If both vanish, write \(k=\kappa c_t\), \(s=\sigma c_m\), and use the unique
decompositions
\[
 \begin{split}
 R&=r_Hc_t+c_mt_H+\rho c_mc_t,\\
 L&=L_H+xc_m,\qquad P=P_H+yc_t,
 \end{split}
\]
with \(r_H\in H_m\), \(t_H\in H_t\),
\(L_H\in V_l\otimes H_m\), \(P_H\in V_l\otimes H_t\).
Then the reduced equation is equivalent to
\begin{equation}\label{atlas:wedge-scalars}
 \kappa L_H+qr_H=0,\qquad
 \sigma P_H+qt_H=0,\qquad
 \kappa x+\sigma y+\rho q=\tau c_l.
\end{equation}
\end{lemma}

\begin{proof}
Write
\(Q=q c_z+c_lu\), \(S=s c_z+c_md\), \(K=k c_z+c_tf\), with
\(u,d,f\in H_z\).  Projection at \(l,m,z\) gives
\(\overline Lc_tf=0\), so \(f=0\); projection at \(l,t,z\) similarly gives
\(d=0\).
We claim that
\begin{equation}\label{atlas:two-term-obstruction}
 Lk+Ps=\tau c_lc_mc_t
\end{equation}
is impossible when both double projections are nonzero.
Its full quotient shows that if one of \(\overline k,\overline s\) is zero,
then both are zero.  Writing \(k=\kappa c_t,s=\sigma c_m\), projection at
\(l,m\) and \(l,t\) forces \(\kappa=\sigma=0\), a contradiction.
If both quotients are nonzero, quotient the whole spaces at \(m,t\) by
\(\C s,\C k\).  The left side vanishes and the pure target survives,
again a contradiction.
Thus \(R\ne0\).  The \(H_z\) component of the original equation is
\(c_luR=0\), giving \(u=0\).  Cancel \(c_z\) to obtain
\eqref{atlas:wedge-reduced}; the same obstruction gives \(q\ne0\).
Full projection of that equation is
\(\overline L\,\overline k+\overline P\,\overline s=0\), proving the
zero/nonzero dichotomy.  Finally the \(H_mc_t,c_mH_t,c_mc_t\) sectors are
exactly \eqref{atlas:wedge-scalars}.
\end{proof}

\begin{proposition}\label{atlas:wedge11}
For \(\rk\overline L=\rk\overline P=1\), the following five families are
disjoint and exhaustive.

\paragraph{I.} Choose \(\eta\ne0\), and take \(L,P\) free with projected
rank one.  Set
\[
 Q=\eta c_lc_z,\qquad R=(\tau/\eta)c_mc_t,\qquad S=K=0.
\]

\paragraph{II.} Choose \(\kappa\ne0\), \(q\in V_l,r\in V_m\) with
\(\overline q,\overline r\ne0\), and take \(P\) free with projected rank one.
Set
\[
 L=(\tau c_lc_m-qr)/\kappa,\quad Q=q c_z,\quad R=r c_t,
 \quad S=0,\quad K=\kappa c_tc_z.
\]

\paragraph{III.} Choose \(\sigma\ne0\), \(q\in V_l,r\in V_t\) with
\(\overline q,\overline r\ne0\), and take \(L\) free with projected rank one.
Set
\[
 P=(\tau c_lc_t-qr)/\sigma,\quad Q=q c_z,\quad R=c_mr,
 \quad S=\sigma c_mc_z,\quad K=0.
\]

\paragraph{IV.} Choose \(\kappa,\sigma\ne0\), \(q\in V_l\) with
\(\overline q\ne0\), nonzero \(A\in H_m,B\in H_t\), arbitrary
\(x\in V_l\) and \(\rho\in\C\).  Put
\[
 y=(\tau c_l-\kappa x-\rho q)/\sigma
\]
and set
\begin{equation}\label{atlas:wedge11-IV}
 \begin{gathered}
 L=qA+xc_m,\qquad P=qB+yc_t,\qquad Q=q c_z,\\
 R=-\kappa Ac_t-\sigma c_mB+\rho c_mc_t,\qquad
 S=\sigma c_mc_z,\qquad K=\kappa c_tc_z.
 \end{gathered}
\end{equation}

\paragraph{V.} Choose \(\eta\ne0\), vectors \(v\in V_l,s\in V_m,k\in V_t\)
with all three quotients nonzero, and arbitrary \(\alpha,\beta\in\C\).
Set
\begin{equation}\label{atlas:wedge11-V}
 \begin{gathered}
 L=vs+\alpha c_lc_m,\qquad P=-vk+\beta c_lc_t,\qquad Q=\eta c_lc_z,\\
 R=(\tau c_mc_t-\alpha c_mk-\beta s c_t)/\eta,\qquad
 S=s c_z,\qquad K=k c_z.
 \end{gathered}
\end{equation}
In IV and V all six whole tensors are nonzero and the two projected
\(l\)-endpoint lines coincide.  Independent \(l\)-endpoint lines occur
only in I--III, with the indicated condition imposed on the free tensors.
\end{proposition}

\begin{proof}
Use Lemma~\ref{atlas:wedge-reduction}.  If \(\overline k=\overline s=0\),
consider \eqref{atlas:wedge-scalars}.
For \(\kappa=\sigma=0\), its first two equations give \(r_H=t_H=0\), and
\(\rho q=\tau c_l\) gives I.
For \(\kappa\ne0,\sigma=0\), we have \(t_H=0\) and
\(\kappa L=\tau c_lc_m-q(r_H+\rho c_m)\).  The nonzero projected rank
forces \(\overline q,r_H\ne0\), giving II.  The opposite scalar choice
gives III.  If both scalars are nonzero, set
\(A=-r_H/\kappa\), \(B=-t_H/\sigma\).  Then the first two equations give
\(L_H=qA,P_H=qB\), with \(\overline q,A,B\ne0\), and the last equation
determines \(y\).  This is IV.

Suppose now \(\overline s,\overline k\ne0\).  Quotient
\eqref{atlas:wedge-reduced} at \(m,t\) by the actual lines \(\C s,\C k\).
The target survives, so the remaining equality \(qR'=\tau c_l[c_m][c_t]\)
forces \(q=\eta c_l\) with \(\eta\ne0\).
Projection at \(l\) gives
\(L_lk+P_ls=0\).  Quotienting at \(m\) by \(\C s\) and then canceling
the nonzero vectors shows
\(L_l=us,P_l=-uk\), for a nonzero \(u\in H_l\).
Consequently \(L=us+c_lA_m\) and \(P=-uk+c_lB_t\).
After substitution, the remaining equation is
\[
 A_mk+sB_t+\eta R=\tau c_mc_t.
\]
Its double quotient is
\(\overline{A_m}\,\overline k+\overline s\,\overline{B_t}=0\).
Because \(\overline s,\overline k\ne0\), contraction by a functional
nonzero on either factor gives
\[
 A_m=\gamma s+\alpha c_m,\qquad
 B_t=-\gamma k+\beta c_t
\]
for scalars \(\gamma,\alpha,\beta\).  Set \(v=u+\gamma c_l\).
This gives exactly V.

For completeness the converses are whole-tensor identities.  In I only
\(qR\) contributes.  In II and III its non-target part cancels the corresponding
part of \(Lk\) or \(Ps\).  In IV the sum is
\[
 (\kappa x+\sigma y+\rho q)c_mc_t=\tau c_lc_mc_t,
\]
since the terms in \(qAc_t\) and \(qc_mB\) cancel separately.
In V the terms \(vsk\), \(\alpha c_lc_mk\), and \(\beta c_ls c_t\)
each cancel, leaving the same target.  Multiplying by \(c_z\) proves
\eqref{atlas:equation}.  The prescribed projected ranks follow from the
displayed nonzero outer factors.

I--III are distinguished by which of \(S,K\) vanish.  In IV both are
nonzero with zero leaf quotients; in V both leaf quotients are nonzero.
Thus the families are disjoint.  IV and V visibly have common projected
\(l\)-endpoint lines.  Conversely, if the lines of \(\overline L\) and
\(\overline P\) are independent, the equation
\(\overline L\,\overline k+\overline P\,\overline s=0\) forces both
leaf quotients zero, and IV still has coincident lines.  This gives the
last assertion.
\end{proof}

\begin{remark}\label{atlas:wedge-parameters}
With the fixed complements, the parameters in each wedge family are
recoverable from the whole tuple.  First recover \(q,s,k\) by contraction
against the covector taking \(c_z\) to one and vanishing on \(H_z\).
In IV, the unique decomposition of \(R\) gives \(A,B,\rho\), and
contraction of \(L\) at \(c_m\) gives \(x\).  In V, contraction of \(L\)
against any covector vanishing on \(c_m\) and taking \(s\) to one recovers
\(v\); the residual pure tensors recover \(\alpha,\beta\).
The simpler recoveries in I--III follow directly from their formulas.
\end{remark}

\begin{proposition}\label{atlas:wedge21}
For \(\rk\overline L=2\), \(\rk\overline P=1\), and the other four
projections zero, \(K=0\) and \(L\) is free with projected rank two.
The remaining tensors have precisely one of the following forms:
\begin{enumerate}
\item \(S=0\), \(Q=a c_lc_z\), \(R=(\tau/a)c_mc_t\), with \(a\ne0\),
and \(P\) free with projected rank one;
\item for \(\sigma\ne0\), \(u\in V_l,v\in V_t\) with
\(\overline u,\overline v\ne0\),
\[
 S=\sigma c_mc_z,\quad Q=u c_z,\quad R=c_mv,\quad
 P=(\tau c_lc_t-uv)/\sigma.
\]
\end{enumerate}
\end{proposition}

\begin{proof}
In \eqref{atlas:wedge-reduced}, write \(\overline P=xy\) with
\(x,y\ne0\).  Full projection gives
\(\overline L\,\overline k+x\overline s\,y=0\).
A nonzero functional at \(l\) annihilating \(x\) cannot annihilate the
rank-two tensor \(\overline L\); hence \(\overline k=0\), and then
\(\overline s=0\).  In \eqref{atlas:wedge-scalars}, projection of the
first equation at \(l\) is a rank-two tensor multiplied by \(\kappa\)
plus one outer product.  Thus \(\kappa=0\), and then \(r_H=0\) because
\(q\ne0\).  If \(\sigma=0\), the remaining equations give the first
branch.  Otherwise write \(u=q\), \(v=t_H+\rho c_t\); they give the second
branch and its nonzero quotient conditions.  Its whole \(P\) has rank two,
since \((c_l,u)\) and \((c_t,v)\) are independent pairs.  Direct substitution
proves both converses, with no further restriction on \(L\).
\end{proof}

\begin{proposition}\label{atlas:wedge22}
For \(\rk\overline L=\rk\overline P=2\), and the other four projections
zero, all solutions are
\[
 S=K=0,\qquad Q=a c_lc_z,\qquad R=(\tau/a)c_mc_t,\qquad a\ne0,
\]
with \(L,P\) free with projected rank two.
\end{proposition}

\begin{proof}
Full projection of \eqref{atlas:wedge-reduced} gives
\(\overline L\,\overline k+\overline P\,\overline s=0\).
If \(\overline k\ne0\), contract at \(t\) with a nonzero functional
annihilating it.  Rank two of \(\overline P\) forces \(\overline s=0\),
and the original equality then forces \(\overline k=0\), a contradiction.
Thus both leaf quotients vanish.  Projecting the first two equations of
\eqref{atlas:wedge-scalars} at \(l\), rank two versus rank at most one
gives \(\kappa=\sigma=0\).  The remaining equations give the stated pure
\(Q,R\).  Their product is the target, proving the converse for arbitrary
whole lifts of \(L,P\).
\end{proof}

\subsection{The two triangle types}

The triangle is on \(l,m,t\), with \(z\) its isolated projected vertex.
The unique pure spoke from Lemma~\ref{support:triangle} distinguishes one
of these three vertices.  In the next statement call that vertex \(t\).
Choosing each of the three possible vertices in turn covers all labelled
triangle solutions.

\begin{proposition}\label{atlas:triangle}
Choose \(\kappa\ne0\), \(q,v\in V_l\), \(s,w\in V_m\), and
\(h\in H_t\), with \(\overline q,\overline s,h\ne0\), and set
\begin{equation}\label{atlas:triangle-form}
 \begin{gathered}
 Q=q c_z,\qquad S=s c_z,\qquad K=\kappa c_tc_z,\\
 P=qh+vc_t,\qquad R=-sh+wc_t,\qquad
 L=(\tau c_lc_m-vs-qw)/\kappa.
 \end{gathered}
\end{equation}
Extend \(\overline q\) to a quotient basis \((\overline q,e)\) at \(l\),
and \(\overline s\) to a quotient basis \((\overline s,f)\) at \(m\).
Write
\[
 \overline v=v_1\overline q+v_2e,\qquad
 \overline w=w_1\overline s+w_2f,\qquad d=v_1+w_1.
\]
The exact additional restrictions for the two triangle types are
\begin{equation}\label{atlas:triangle-test}
 \begin{array}{ll}
 (1,1,1):&v_2w_2=0,\quad(d,v_2,w_2)\ne(0,0,0),\\
 (2,1,1):&v_2w_2\ne0.
 \end{array}
\end{equation}
In the second row \(L\) is the rank-two projected edge.
These formulas exhaust both types and give their full converses.
For the fixed complement \(H_t\), their parameters are unique once the
pure-spoke vertex and the ordering of the other two vertices are fixed.
\end{proposition}

\begin{proof}
Apply Lemma~\ref{support:triangle}, with its pure-spoke vertex taken as \(t\).
It gives the whole forms \eqref{atlas:triangle-form} with a lift of a
nonzero quotient vector in the \(t\) slot.  Changing that lift by
\(h\mapsto h+\lambda c_t\) changes the remaining vectors by
\[
 v\mapsto v-\lambda q,\qquad w\mapsto w+\lambda s.
\]
These simultaneous changes leave every tensor unchanged, since
\(vs+qw\) is unchanged.  They allow and uniquely determine the choice
\(h\in H_t\).

The two incident projections are
\(\overline P=\overline q\,h\) and
\(\overline R=-\overline s\,h\), both nonzero of rank one.  The unscaled
numerator of the opposite double projection has matrix
\[
 \overline v\,\overline s+\overline q\,\overline w
 \quad\longleftrightarrow\quad
 \begin{pmatrix}d&w_2\\v_2&0\end{pmatrix}.
\]
Its determinant is \(-v_2w_2\).  A nonzero matrix of this size has rank one
exactly when its determinant vanishes, and rank two exactly when it does
not.  This proves \eqref{atlas:triangle-test}, including the branches
\(v_2=0\) and \(w_2=0\) and their intersection; the excluded zero triple
is exactly a vanished opposite projection.

Conversely the three products from \eqref{atlas:triangle-form} are
\[
\begin{aligned}
 LK&=\tau c_lc_mc_tc_z-vsc_tc_z-qwc_tc_z,\\
 PS&=qshc_z+vsc_tc_z,\\
 QR&=-qshc_z+qwc_tc_z.
\end{aligned}
\]
Their sum is the required whole target, with all scalar components retained.
The rank test gives the desired support exactly.
Finally \(Q,S,K\) recover \(q,s,\kappa\) by contraction at \(c_z\).
The partial \(t\)-projection of \(P\) recovers the quotient of \(h\),
and the fixed complement recovers \(h\) itself.  The residual tensors in
\(P,R\) then recover \(v,w\).  This proves uniqueness.
\end{proof}

\begin{remark}\label{atlas:triangle-ranks}
All six tensors and all three products are nonzero in either triangle type.
For type \((2,1,1)\), the pairs \((\overline q,\overline v)\) and
\((\overline s,\overline w)\) are independent.  In the bases
\((c_l,q,v)\), \((c_m,s,w)\), the matrix of \(L\) is
\[
 \kappa^{-1}
 \begin{pmatrix}\tau&0&0\\0&0&-1\\0&-1&0\end{pmatrix}.
\]
Thus its whole rank is three, while \(P,R\) have whole rank two and the
spokes have rank one.  For type \((1,1,1)\), either
\(\Span(c_l,q,v)\) or \(\Span(c_m,s,w)\) has dimension at most two,
so all three triangle tensors have rank at most two.  Vanishing individual
coordinates or further whole-rank drops require no additional chart.
\end{remark}

\subsection{The four-cycle}

\begin{proposition}\label{atlas:cycle}
Suppose \(\overline L=\overline K=0\), and
\(\overline P,\overline Q,\overline R,\overline S\) are nonzero of rank one.
All solutions are exactly
\begin{equation}\label{atlas:cycle-form}
 \begin{gathered}
 L=\gamma c_lc_m,\qquad K=(\tau/\gamma)c_tc_z,\qquad\gamma\ne0,\\
 P=u_lu_t,\qquad Q=u_lu_z,\qquad
 R=u_mu_t,\qquad S=-u_mu_z,
 \end{gathered}
\end{equation}
where \(u_i\in V_i\) and \(\overline{u_i}\ne0\) for all four sites.
All six whole tensors have rank one, and
\(PS+QR=0\), \(LK=\tau c_lc_mc_tc_z\).
\end{proposition}

\begin{proof}
Full projection gives
\(\overline P\,\overline S=-\overline Q\,\overline R\ne0\).
Equality of nonzero simple four-tensors gives equal endpoint lines at
every site.  The scalar factors can be chosen consistently: first factor
\(\overline P=a_la_t\), then \(\overline Q=a_la_z\), then
\(\overline R=a_ma_t\); the equality forces
\(\overline S=-a_ma_z\).  Here each \(a_i\) is a nonzero vector in \(H_i\).
Retain all whole lifts by writing
\[
\begin{aligned}
 P&=a_la_t+p_lc_t+c_lp_t+p_0c_lc_t,\\
 Q&=a_la_z+q_lc_z+c_lq_z+q_0c_lc_z,\\
 R&=a_ma_t+r_mc_t+c_mr_t+r_0c_mc_t,\\
 S&=-a_ma_z+s_mc_z+c_ms_z+s_0c_mc_z,
\end{aligned}
\]
with the subscripted vector at site \(i\) in \(H_i\).
Since \(\overline L=\overline K=0\), their product contributes no
three-\(H\) component.  The four three-\(H\) components of the full equation
therefore give
\[
 a_ls_m+q_la_m=0,\quad -p_la_m+a_lr_m=0,\quad
 a_ts_z+r_ta_z=0,\quad -p_ta_z+a_tq_z=0.
\]
For nonzero \(a,b\), an equality \(ax+yb=0\) forces
\(x=\lambda b,y=-\lambda a\), including \(\lambda=0\): quotienting by
either fixed line and then canceling the other factor proves this.
Consequently there are scalars \(\alpha,\beta,\delta,\epsilon\) with
\[
\begin{gathered}
 p_l=\alpha a_l,\quad r_m=\alpha a_m,\quad
 q_l=-\beta a_l,\quad s_m=\beta a_m,\\
 p_t=\delta a_t,\quad q_z=\delta a_z,\quad
 r_t=\epsilon a_t,\quad s_z=-\epsilon a_z.
\end{gathered}
\]
Define whole vectors
\[
 x_l=a_l+\delta c_l,\quad x_m=a_m+\epsilon c_m,\quad
 x_t=a_t+\alpha c_t,\quad x_z=a_z-\beta c_z.
\]
Expansion gives
\begin{equation}\label{atlas:cycle-residual}
 \begin{gathered}
 P=x_lx_t+A c_lc_t,\qquad Q=x_lx_z+B c_lc_z,\\
 R=x_mx_t+C c_mc_t,\qquad S=-x_mx_z+D c_mc_z,
 \end{gathered}
\end{equation}
where
\(A=p_0-\delta\alpha\), \(B=q_0+\delta\beta\),
\(C=r_0-\epsilon\alpha\), \(D=s_0-\epsilon\beta\).
All \(\overline{x_i}\) are nonzero.

Choose complements to \(\C c_i\) containing \(x_i\); this is merely a
choice of basis for the following coefficient comparison, and does not
change the tensors or \(\tau\).  Write the still unrestricted missing edges
in these complements as
\[
 L=\xi_lc_m+c_l\eta_m+\ell c_lc_m,\qquad
 K=\rho_tc_z+c_t\sigma_z+\kappa c_tc_z.
\]
The four cross-pair components, the four single-site components, and the
constant component of \eqref{atlas:equation} are respectively
\begin{equation}\label{atlas:cycle-coefficients}
\begin{gathered}
 \xi_l\rho_t=-D x_lx_t,\qquad
 \xi_l\sigma_z=-C x_lx_z,\\
 \eta_m\rho_t=-B x_mx_t,\qquad
 \eta_m\sigma_z=A x_mx_z,\\
 \kappa\xi_l=\kappa\eta_m=\ell\rho_t=\ell\sigma_z=0,\qquad
 \ell\kappa+AD+BC=\tau.
\end{gathered}
\end{equation}
These are all remaining coefficients: the other two pair components are
identically zero, the three-site components have just been eliminated,
and the four-site components cancel in \eqref{atlas:cycle-residual}.

Multiply the first and fourth pair equations in
\eqref{atlas:cycle-coefficients}, and then the second and third, restoring
both to the same site order.  Their left sides are the identical tensor
\(\xi_l\eta_m\rho_t\sigma_z\).  Since \(x_lx_mx_tx_z\ne0\), their right
sides give \(-DA=CB\), so \(AD+BC=0\).
This argument divides by none of the unknown vectors or residual scalars.
The final equation now gives \(\ell\kappa=\tau\ne0\).
The single-site equations force \(\xi_l=\eta_m=\rho_t=\sigma_z=0\),
and the pair equations force \(A=B=C=D=0\).
Taking \(\gamma=\ell\) and \(u_i=x_i\) yields
\eqref{atlas:cycle-form} in the original spaces.

Conversely let the parameters in \eqref{atlas:cycle-form} be given, and put
\(U=u_lu_mu_tu_z\).  Then \(LK=\tau c_lc_mc_tc_z\), \(PS=-U\), and
\(QR=U\).  The double projections have exactly the prescribed ranks,
and every whole factor is nonzero of rank one.  This proves the full
converse, including arbitrary \(c\)-components in the four vectors.
\end{proof}

The propositions describe whole tensors, rather than only quotient data.
Every displayed parametrization therefore includes all permitted cancellations
and whole-rank drops.  The list is finite at the level of support and rank
types; Proposition~\ref{atlas:empty} uses geometric parameters, and no
classification into finitely many orbits is asserted.


\section{Dimensions of the strata}\label{geom:section}

We use algebraic dimension for constructible strata, equivalently the
dimension of their closures. A whole block with zero double projection
belongs to a five-dimensional vector space. The loci of blocks with
projected ranks one and two have dimensions \(8\) and \(9\), respectively:
the corresponding loci in \(\operatorname{Mat}_{2\times2}(\C)\) have
dimensions \(3\) and \(4\), and the five target-containing entries are free.

\subsection{The empty projected support}

The empty-support chart identifies a point with a vector
\(a\in E=\C^6\) and four maps \(B_i:H_i^*\to S_i\), where \(S_i\)
is the coordinate star at \(i\), satisfying
\begin{equation}\label{geom:empty}
 q(a)=\tau,\qquad U=\sum_i\im B_i\subset a^\perp,\qquad q|_U=0,
 \qquad q(x)=x_{12}x_{34}+x_{13}x_{24}+x_{14}x_{23}.
\end{equation}
Here perpendicularity uses the polar form of \(q\).
The subspace \(a^\perp\) is nondegenerate of dimension five, so
\(r=\dim U\le2\).

\begin{proposition}\label{geom:empty-dimension}
The stratum with all six double projections zero has dimension \(11\).
\end{proposition}
\begin{proof}
The actual maps determine \(U\). On fixed-rank Grassmannian charts,
recovery of their image is regular. We may therefore count the incidence
data \((U,a,(B_i))\), imposing the open condition that the images span \(U\),
without adding an artificial subspace parameter.

Identify \(E\) with \(\bigwedge^2\C^4\), changing the sign of the \(13\)
coordinate. Then \(q\) becomes the Pfaffian, and \(S_i=e_i\wedge\C^4\).
A nonzero isotropic bivector is decomposable. An isotropic two-plane \(U\)
is consequently of the form
\[
 U=\ell\wedge H,\qquad \ell\subset H\subset\C^4,\quad
 \dim\ell=1,\quad\dim H=3,
\]
for a unique flag. Indeed two independent decomposable generators have
wedge zero, so their supporting two-planes meet in a line and span a
three-space. This also proves the converse.

First suppose \(r=2\) and \(\ell=\C e_i\).
Then \(U\subset S_i\). It contains \(m=0,1,\) or \(2\) of the three
coordinate axes of \(S_i\), and these alternatives have parameter
dimensions \(2-m\). For \(j\ne i\), \(U\cap S_j\) is a line precisely
when the corresponding coordinate edge lies in \(U\). Hence
\(\sum_j\dim(U\cap S_j)=2+m\).
For any fixed isotropic two-plane, \(U^\perp\) has dimension four,
radical \(U\), and a nondegenerate two-dimensional quotient.
The set \(q(a)=\tau\) in \(U^\perp\) is thus nonempty of dimension three.
The maps \(B_j\) contribute \(2(2+m)\) parameters. The total dimension is
\[
 (2-m)+3+2(2+m)=9+m\le11.
\]
The spanning condition is a nonempty open condition.

If \(\ell\) is not a coordinate line, then \(U\) is contained in no star,
and \(U\cap S_i\) has dimension one exactly when \(e_i\in H\).
Let \(k\) be the number of coordinate axes contained in \(H\).
Spanning \(U\) by these intersections requires \(k\ge2\), and \(k\le3\).
For \(k=2\), \(H\) varies in an open subset of \(\mathbf P^1\);
\(\ell\) varies in \(\mathbf P(H)\) outside the plane spanned by those
two axes, so the flag parameter dimension is three. The maps contribute
four and \(a\) contributes three, for a total of ten.
For \(k=3\), \(H\) is a coordinate three-space and \(\ell\) varies in
its projective plane, excluding the coordinate points. The flag dimension
is two, the maps contribute six, and \(a\) contributes three, for a total
of eleven.
Here the three nonzero intersections are the lines \(\ell\wedge e_i\) for
the three axes \(\C e_i\subset H\); wedging with \(\ell\) identifies
\(H/\ell\) with \(U\), and the images of these axes span \(H/\ell\), so
the three lines span \(U\). If \(\ell\) lies in the plane of two of these
axes, their two lines coincide and the third differs; each \(U\cap S_i\)
is still a line, so the maps still contribute six and spanning is still
a nonempty open condition.
These are all the two-plane cases.

If \(r=1\), \(U\) must lie in at least one star. A non-coordinate line
lies in exactly one star and has two line parameters and two map
parameters. A coordinate edge line has no line parameter and four map
parameters, from its two incident stars. For either type of isotropic
line, \(q(a)=\tau\) in \(U^\perp\) has dimension four. Thus \(r=1\)
contributes dimension eight. For \(r=0\) the dimension is five.

The incidence alternatives are finite and retain all intersection jumps.
They give an upper bound of eleven. Equality is also visible in
\eqref{main:large} by requiring both free blocks to have zero double
projection: this is \(\A^{10}\times\Gm\), of dimension eleven.
\end{proof}

\subsection{The remaining charts}\label{geom:counts}

The following counts use fixed complements \(H_i\) and whole-vector
parameters. They do not count arbitrary changes of basis as extra
parameters. In every chart whose dimension is asserted, the parameters
are recoverable regularly on finitely many nonzero-coordinate charts.
This recovery is given for the wedges of ranks \((1,1)\) in
Remark~\ref{atlas:wedge-parameters} and for the triangles in
Proposition~\ref{atlas:triangle}; for the single-edge charts, the other
wedges and the rows of the crossing table it is written below, where each
is counted; for the four-cycle, the end of this subsection identifies the
redundancy exactly.

For a single rank-one projected edge \(L\), write
\(K=\kappa c_tc_z\). When \(\kappa\ne0\) and both crossing products are
nonzero, the four vectors \(p,q,r,s\) contribute twelve parameters and
\(\kappa\) contributes one. The condition
\[
 \rk\bigl(\bar p\bar s+\bar q\bar r\bigr)=1
\]
cuts a genuine hypersurface: the determinant is, up to sign, the product
of the two determinants of the endpoint vector pairs. The nonzero-matrix
condition removes neither entire hypersurface component. This chart has
dimension twelve. With one crossing product zero, there are \(\kappa\),
two whole vectors, and a free zero-projection partner, again giving
\(1+3+3+5=12\).
In these \(\kappa\ne0\) charts \(K\) gives \(\kappa\) as its \(c_tc_z\)
coefficient, contraction at \(c_t\) or \(c_z\) recovers each named vector
from its block among \(P,Q,R,S\), and a free partner is itself a block.
When \(\kappa=0\), the free rank-one-projected \(L\) contributes eight.
A zero crossing product leaves the scalar of the two pure blocks whose
product is the target, and a free zero-projection block, giving
\(8+1+5=14\). If both products are nonzero,
each exceptional-site chart has two nonzero scalars and one free whole
vector, with the other vector solved by the target equation; its
dimension is \(8+5=13\).
For \(\kappa=0\), \(L\) and any free block are blocks of the tuple, and
the scalar of a zero crossing product is the coefficient of the pure
block \(Q\) in the first case of Lemma~\ref{atlas:crossing} and of \(P\)
in the second. In each row of the crossing table the two scalars are the
coefficients of its two pure blocks, and contraction at the \(c\) factor
shown there recovers \(u\) from \(P\) or \(S\).
The single rank-one edge stratum therefore has dimension fourteen.

For a single rank-two projected edge, the \(\kappa\ne0\) chart has
dimension thirteen, since its endpoint independence conditions are open.
For \(\kappa=0\), the free block has dimension nine, so the two preceding
crossing alternatives give dimensions fifteen and fourteen.
This stratum has dimension fifteen.
In its charts the parameters are recovered as for rank one.

For a rank-\((1,1)\) wedge, the five families in the atlas have dimensions
\begin{equation}\label{geom:wedge-counts}
 17,\quad15,\quad15,\quad13,\quad12,
\end{equation}
respectively. In the first family these are two free blocks of dimension
eight and one nonzero scalar. In the second and third they are a free
block of dimension eight, two whole vectors and a nonzero scalar.
In the fourth, the free parameters are two scalars, a whole vector \(q\),
two vectors \(A\in H_m,B\in H_t\), a whole vector \(x\), and one scalar
\(\rho\); their dimensions sum to \(2+3+2+2+3+1=13\), and \(y\) is
determined. In the fifth they are \(\eta,v,s,k,\alpha,\beta\), of total
dimension \(1+3+3+3+1+1=12\).
All stated nonvanishing conditions are open, and the fixed-coordinate
recovery in the atlas removes any ambiguity.

A rank-\((2,1)\) wedge has two charts. Their dimensions are
\(9+8+1=18\) and \(9+1+3+3=16\).
A rank-\((2,2)\) wedge has dimension \(9+9+1=19\).
These counts include the arbitrary whole lifts of its free blocks.
In these charts the free blocks are blocks of the tuple. The scalar is
the \(c_lc_z\) coefficient of \(Q\), except in the second \((2,1)\)
chart, where it is the \(c_mc_z\) coefficient of \(S\) and contraction
of \(Q\) at \(c_z\) and of \(R\) at \(c_m\) recovers the two vectors.

For triangles use, for example, the pure spoke \(K=\kappa c_tc_z\).
The common chart is
\begin{equation}\label{geom:triangle-chart}
 \begin{gathered}
 Q=q c_z,\quad S=s c_z,\quad K=\kappa c_tc_z,\\
 P=qh+v c_t,\quad R=-sh+w c_t,\quad
 L=\kappa^{-1}(\tau c_lc_m-vs-qw).
 \end{gathered}
\end{equation}
Require \(h\in H_t\setminus\{0\}\); this removes the entire additive
ambiguity in choosing its whole lift. The parameters
\(q,s,v,w,h,\kappa\) have total dimension fifteen.
For triangle ranks \(211\), the two quotient endpoint pairs
\((\bar q,\bar v)\) and \((\bar s,\bar w)\) are independent; these
are open nonempty conditions. For ranks \(111\), the opposite projected
matrix is nonzero and its determinant vanishes. Its determinant factors
as the product of those two endpoint determinants, giving a union of
hypersurfaces in this parameter domain. Its dimension is fourteen.
There are only three possible pure-spoke branches for a fixed triangle.

Finally the four-cycle chart has a nonzero scalar and four whole vectors
outside their target lines, giving thirteen parameters. The entire
redundancy is
\[
 (u_l,u_m,u_t,u_z)\longmapsto
 (a u_l,a u_m,a^{-1}u_t,a^{-1}u_z),\qquad a\in\Gm.
\]
Equality of the three products \(u_lu_t,u_lu_z,u_mu_t\) forces exactly
this relation; the fourth then agrees. Normalizing a nonzero quotient
coordinate of \(u_l\) gives finite parameter charts of dimension twelve.

Together with Proposition~\ref{geom:empty-dimension}, these calculations
prove the full-stratum dimension column of Table~\ref{main:table}.

\section{Maximal components and nonzero summands}

\begin{proof}[Proof of Theorem~\ref{main:geometry}(1)]
For each centred wedge let \(X_{i;jk}\) be the set
\eqref{main:large}. It is closed in \(\Ftau\): impose the two zero-block
conditions and the two pure-line conditions. The fibre equation then
requires the two pure scalar coefficients to satisfy \(qr=\tau\).
Consequently
\[
 X_{i;jk}\simeq\A^{18}\times
 \operatorname{Spec}\C[q,r]/(qr-\tau)
 \simeq\A^{18}\times\Gm.
\]
Its rank-\((2,2)\) projected subset is nonempty and dense, so its closure
is exactly \(X_{i;jk}\). There are \(4\binom32=12\) centred wedges.
Their varieties are distinct because their dense projected supports differ.

The fibre is a finite union of its labelled strata.
All strata other than rank-\((2,2)\) wedges have dimension at most eighteen.
Taking closures gives a finite closed cover by the twelve \(X_{i;jk}\)
and sets of dimension at most eighteen. Hence \(\dim\Ftau=19\).
An irreducible component of dimension nineteen must be contained in one
of these twelve closed irreducible sets, and maximality makes it equal
to that set. Conversely each \(X_{i;jk}\) is such a component, since any
component containing it has dimension nineteen and must be one of the
twelve. This proves both assertions.
\end{proof}

\begin{proof}[Proof of Theorem~\ref{main:geometry}(2)]
Intersect every chart with the condition that all six whole blocks are
nonzero. For a single projected edge only the \(\kappa\ne0\) chart with
both crossing products nonzero survives, giving dimensions twelve and
thirteen for projected ranks one and two.
For a rank-\((1,1)\) wedge only its fourth and fifth families survive,
of dimensions thirteen and twelve. Wedges of ranks \(21\) and \(22\)
force a whole block to vanish. Both triangle types and the four-cycle
are entirely nonzero-block charts. The empty support has dimension at
most eleven. This proves the last column of Table~\ref{main:table}
and gives \(\dim\Fact=15\).

The dimension-fifteen strata are precisely the rank-\((2,1,1)\) triangles.
There are twelve labelled choices: an isolated projected vertex and a
high edge among the other three vertices.
For any fixed choice, the chart \eqref{geom:triangle-chart}, with the
fixed-complement convention, has parameter domain
\[
 \kappa\in\Gm,\quad h\in H_t\setminus\{0\},\quad
 (\bar q,\bar v),(\bar s,\bar w)\text{ independent}.
\]
It is a nonempty open subset of affine fifteen-space. The recovery of
its actual parameters is regular on nonzero-coordinate charts. Hence
this stratum is irreducible of dimension fifteen.

Its closure within \(\Fact\) is therefore irreducible of dimension fifteen.
These twelve closures are distinct. Different isolated vertices impose
different closed zero-projection equations. For a fixed isolated vertex,
different high edges are separated by the closed rank-at-most-one
conditions on the two low edges. Each other triangle stratum violates
one of the relevant closed conditions.
All remaining active strata have dimension at most fourteen.
The same finite-closed-cover argument as above now proves that the
twelve closures are exactly the components of maximal dimension in
\(\Fact\).
\end{proof}

\begin{proof}[Proof of Theorem~\ref{main:geometry}(3)]
For nonzero vector spaces \(U,V\), the factorization map
\[
 (U\setminus0)\times(V\setminus0)\longrightarrow
 \{a\otimes b:a,b\ne0\}
\]
has fibre \(\Gm\), acting by \((a,b)\mapsto(sa,s^{-1}b)\).
It is Zariski locally trivial. In chosen bases, on the chart \(T_{pq}\ne0\)
one normalized factorization is
\[
 a_p=1,\qquad b_j=T_{pj},\qquad a_i=T_{iq}/T_{pq}.
\]
The rank-one equations imply \(T_{ij}=a_i b_j\).

Apply this independently to each of the three paired partitions and
restrict to \(A+B+C=\tau c_1c_2c_3c_4\).
The result is a surjective locally trivial \((\Gm)^3\)-bundle
\(\Fact\to\Ddec\). Dimensions therefore differ by three.
The bundle also induces a bijection on irreducible components.
Indeed its inverse image over an irreducible base component is locally
a product with an irreducible torus; a dense trivializing open set has
irreducible inverse image, whose closure fills that inverse image.
For the last claim, a point of that inverse image lies over a second
trivializing open set; the two sets meet the irreducible base component
in a common dense open subset, and in the product chart over the second
set the inverse image of that subset is dense, so the point lies in the
closure.
Conversely maximal irreducible closed subsets upstairs project densely
onto base components, by the local product description.
Thus all component dimensions decrease by three, proving the claim.
\end{proof}

\begin{remark}
The twelve largest components of \(\Ftau\) are disjoint from \(\Fact\).
The active components of dimension fifteen therefore come from
lower-dimensional components of the full fibre. We do not claim that
every active point belongs to one of the twelve largest active components,
nor identify their boundary intersections.
\end{remark}

\section{Examples and scope}\label{examples:section}

For a rank-\((2,1,1)\) triangle, choose bases \((c_i,a_i,b_i)\) and set
\[
 \begin{gathered}
 Q=a_lc_z,\quad S=a_mc_z,\quad K=c_tc_z,\\
 P=a_la_t+b_lc_t,\qquad
 R=-a_ma_t+b_mc_t,\\
 L=\tau c_lc_m-b_la_m-a_lb_m.
 \end{gathered}
\]
Direct substitution gives \(LK+PS+QR=\tau c_lc_mc_tc_z\).
The whole ranks of \(L,P,Q,R,S,K\) are respectively \(3,2,1,2,1,1\).
The projected triangle has ranks \(2,1,1\); every matching product is
nonzero. The open family \eqref{geom:triangle-chart} explains the
fifteen-dimensional stratum containing such points.

The empty projected support can also exhibit nonzero cancellation.
Use one quotient variable \(x\) at site \(1\) and one \(y\) at site \(2\),
and consider the affine edge forms
\[
 E_{12}=1-2y,\quad E_{13}=1+x,\quad E_{14}=1-x,\quad
 E_{23}=1+y,\quad E_{24}=1+y,\quad E_{34}=1.
\]
They satisfy
\[
 E_{12}E_{34}+E_{13}E_{24}+E_{14}E_{23}=3.
\]
Homogenizing separately at each site with its coordinate \(c_i\) gives
six nonzero whole edge tensors with zero double projections. The
corresponding isotropic image \(U\) has dimension two.
Multiplying all six edges by a scalar \(s\) with \(3s^2=\tau\) gives the
prescribed target. This example illustrates why projected support
alone does not determine whole-block support.

The assumptions \(\tau\ne0\) and \(\dim V_i=3\) are used explicitly.
The nonzero target is essential in the contradiction that proves
Lemma~\ref{support:degree} and in the reductions to a pure spoke. The
quotient dimension two bounds the projected ranks and enters the incidence
count. We make no assertion
for the zero fibre, for other local dimensions, or for scheme
multiplicities. The lower-dimensional components and the incidence
relations between the listed strata remain further questions.

\section*{Proof artifacts and AI use}
The author, who publishes under the pseudonym Ignoramus, directed this
project and takes responsibility for the paper. The author used OpenAI
Codex to produce the results, proofs, checking program and text of this
paper, and Anthropic Claude to check and revise it as described below,
and has not verified the mathematics; no human has checked the proofs. No
human peer review or formal proof-assistant verification is claimed, and
no AI system is listed as an author. The attribution to OpenAI Codex, of
this work and of the Codex checks described below, rests on the project's
records and the author's account; which model versions were used at each
step has not been verified.

All checks described in this section were made by AI agents. The
classifications of the single edge of projected rank two, of the wedge of
projected ranks \((2,1)\) and of the empty projected support,
Lemma~\ref{support:triangle} in its present form, the assembly into nine
types, and every dimension and component count,
Theorem~\ref{main:geometry} included, were first obtained in the review
that led to this paper, shortly before it was drafted; the other results
go back to earlier stages of the project, some of them in a weaker form.
In that review, before drafting, agents of the Codex workflow checked the
arguments and compared the results with published work. Two further
agents of that workflow then reviewed the manuscript: one reported that
it had reconstructed the proofs analytically, claiming no independent
machine computation, and that it had found no error or gap requiring
repair; the other checked the claims of novelty, the citations and the
wording. Both had seen the workflow's
earlier assessments of the results, so neither review was blind.

Anthropic Claude prepared this revision at the author's instruction. It
wrote this section and the abstract. In Section~1 it wrote the paragraph
that defines pure tensors, the definitions of terms, the sentence that
names the components of dimension fifteen, the sentence on searches for
earlier statements, and the sentences that cite the author's earlier
papers and~\cite{kgz}; it added the reference~\cite{naslund}, made the
pointers to the literature exact and reworded the comparison with the
note~\cite{axis}. It also changed wording in other places and renamed
letters that carried two meanings in the proofs of
Lemmas~\ref{support:degree} and~\ref{support:triangle} and in the
treatment of the empty projected support, which changed letters in six
displayed formulas. It added short sentences that write out steps the
earlier text left to the reader: the uniqueness in
Lemma~\ref{support:triangle}, an instance of each of the nine types, one
point each in the proofs of Propositions~\ref{atlas:empty}
and~\ref{geom:empty-dimension} and of Theorem~\ref{main:geometry}, and,
in Section~\ref{geom:counts}, how the parameters are recovered in the
single-edge charts, in the rows of the crossing table and in the wedges
of ranks \((2,1)\) and \((2,2)\). Exact programs, not posted with the
paper, tested these sentences on examples; for the step in the proof of
Theorem~\ref{main:geometry}(3) they tested only the two facts it uses.
It changed no statement of a theorem, proposition, lemma or remark, and
it removed or altered no step of the earlier proofs.

Claude agents also checked the paper. One read every line of the
unrevised text and found no error and no gap affecting a theorem. One
compared the paper with the author's four earlier papers and found none
of its results there. A fact-check compared the citations then in the
text with the cited works and led to corrections. The Claude session that
coordinated this revision compared its own sentences citing the author's
earlier papers with those papers, and an earlier draft of this section
with the project's records; a further agent audited every verdict, and
this section was corrected. Three agents, each working alone, then read
the whole text as critics, before the corrections that followed and
before the short added sentences described above were written; each
checked every numbered statement by hand and tested the
parametrizations, the dimensions and the list of nine types with programs
of its own. None found a mathematical error; each found faults of
wording, credit and definition, most of which were then corrected. The
corrections and the added sentences were then compared line by line with
the unrevised text by one further Claude agent that did not write them;
it checked the added sentences by hand and with programs of its own, and
the references added after the fact-check against their sources.
Eight of the ten corrections it proposed were then made, in this section
and in four sentences elsewhere. Of the other two, the one on the
attribution to OpenAI Codex was left to the author, and the other was set
aside on the record of an earlier check. None of the checks named in this
section is journal peer review, a formal verification or a certificate of
originality or priority.

No step of a proof depends on machine computation. The program
\texttt{check\_examples.py}, posted with the paper together with its
saved results and a README, checks 20 examples with rational entries,
covering all nine projected support and rank types, each under all 24
permutations of the four sites, in exact rational arithmetic without
floating point. For each example and permutation it checks all 81
coefficients of \(LK+PS+QR=\tau c_lc_mc_tc_z\), the whole and projected
ranks of the six blocks, the projected support type, and which blocks and
which matching products are zero. For the empty-support examples it also
checks, in the notation of Proposition~\ref{atlas:empty}, that
\(q(a)=\tau\) and that \(\sum_i\im B_i\) is totally isotropic, orthogonal
to \(a\) and of dimension two. The examples include the two of the
preceding section. All but one use \(\tau=3\); one uses \(\tau=12\). These
are finite checks of examples. They do not prove that the nine types are
exhaustive or that the parametrizations are complete, and they do not
check the dimension or component counts.

\begin{thebibliography}{99}
\raggedright
\bibitem{ballico}
E.~Ballico, \emph{Joins, Secant Varieties and Their Associated Grassmannians},
Mathematics \textbf{12} (2024), 1274.
\href{https://doi.org/10.3390/math12091274}{doi:10.3390/math12091274}.

\bibitem{ghl}
F.~Gesmundo, Y.~I.~Han and B.~Lovitz,
\emph{Linear preservers of secant varieties and other varieties of tensors},
Journal of Symbolic Computation \textbf{131} (2025), 102449;
\href{https://doi.org/10.1016/j.jsc.2025.102449}{doi:10.1016/j.jsc.2025.102449}.
arXiv:2407.16767.

\bibitem{karam}
T.~Karam, \emph{Small sunflowers and the structure of slice rank decompositions},
Journal of the London Mathematical Society \textbf{113} (2026), e70485.
\href{https://doi.org/10.1112/jlms.70485}{doi:10.1112/jlms.70485}.
arXiv:2308.07101.

\bibitem{lovitzpetrov}
B.~Lovitz and F.~Petrov,
\emph{A generalization of Kruskal's theorem on tensor decomposition},
Forum of Mathematics, Sigma \textbf{11} (2023), e27, 1--40.
\href{https://doi.org/10.1017/fms.2023.20}{doi:10.1017/fms.2023.20}.

\bibitem{axis}
djh58 (repository account), \emph{The general axis bound \(D\le N-2\)},
provisional, unreviewed repository note,
\texttt{krenn-gu-n6d4/proof/general\_bound.md}.
\url{https://github.com/djh58/krenn-gu-n6d4/blob/main/proof/general_bound.md}
(accessed 10 November 2026).

\bibitem{ekm}
R.~Elman, N.~Karpenko and A.~Merkurjev,
\emph{The Algebraic and Geometric Theory of Quadratic Forms},
American Mathematical Society Colloquium Publications, vol.~56,
American Mathematical Society, 2008.

\bibitem{ign-geometry}
Ignoramus, \emph{Dimensions and generic recovery of matching tensors},
preprint, Zenodo, 2026.
\href{https://doi.org/10.5281/zenodo.24114656}{doi:10.5281/zenodo.24114656}.

\bibitem{ign-relations}
Ignoramus, \emph{Complementary tensor relations and bipartite matching
kernels}, version~3.3, preprint, Zenodo, 2026.
\href{https://doi.org/10.5281/zenodo.24115283}{doi:10.5281/zenodo.24115283}.

\bibitem{ign-boundary}
Ignoramus, \emph{Internal-vertex complexity of four-boundary matching
signatures}, preprint, Zenodo, 2026.
\href{https://doi.org/10.5281/zenodo.24116409}{doi:10.5281/zenodo.24116409}.

\bibitem{kgz}
M.~Krenn, X.~Gu and A.~Zeilinger,
\emph{Quantum experiments and graphs: Multiparty states as coherent
superpositions of perfect matchings},
Physical Review Letters \textbf{119} (2017), 240403.
\href{https://doi.org/10.1103/PhysRevLett.119.240403}{doi:10.1103/PhysRevLett.119.240403}.

\bibitem{naslund}
E.~Naslund,
\emph{The partition rank of a tensor and \(k\)-right corners in \(\mathbb F_q^n\)},
Journal of Combinatorial Theory, Series A \textbf{174} (2020), 105190.
\href{https://doi.org/10.1016/j.jcta.2019.105190}{doi:10.1016/j.jcta.2019.105190}.
\end{thebibliography}
\end{document}

